\chapter{Extracorporeal Membrane Oxygenation (ECMO)}

\section*{Overview}
Extracorporeal membrane oxygenation (ECMO) is a form of mechanical life support that temporarily replaces the function of the heart or lungs using an external pump and oxygenator. Blood is drained from the patient, oxygenated outside the body, and returned to the circulation.

\section*{Alternative Names}
ECMO, extracorporeal life support (ECLS), venovenous ECMO (VV-ECMO), venoarterial ECMO (VA-ECMO), heart-lung bypass

\section*{Principle}
A membrane oxygenator adds oxygen and removes carbon dioxide from blood while a mechanical pump maintains flow, allowing the native heart and lungs time to recover. Venovenous (VV) ECMO provides isolated respiratory support, while venoarterial (VA) ECMO supports both cardiac and pulmonary function. Cannula placement and circuit configuration determine the mode of support.

\section*{History}
The first successful clinical use of ECMO occurred in 1972, with early trials in neonatal respiratory failure showing benefit. Its use expanded substantially following the 2009 H1N1 influenza pandemic and the COVID-19 pandemic.

\section*{Why It Is Done}
Initiated for severe, potentially reversible respiratory or cardiac failure refractory to maximal conventional therapy, including hypoxemic or hypercapnic respiratory failure, cardiogenic shock, and cardiac arrest. It also supports patients as a bridge to transplant, to ventricular assist device placement, or to recovery.

\section*{Is This a Primary or Secondary Test or Procedure}
A rescue therapy used when primary treatments fail; it is supportive, not curative, and always temporary.

\section*{How It Is Performed}
Cannulas are inserted percutaneously or surgically into large veins (VV-ECMO) or a vein and artery (VA-ECMO). Blood flows through the circuit to the pump and oxygenator and returns warmed to the patient, while anticoagulation is infused to prevent thrombosis. A perfusionist and critical care team manage the circuit continuously.

\section*{Preparation}
Echocardiography or bedside ultrasound to guide cannulation and rule out right heart thrombus. Confirmation of type and crossmatch compatibility for circuit priming, and placement of invasive monitoring lines.

\section*{Contraindications and Precautions}
Contraindications include irreversible end-organ failure, advanced age with prohibitive comorbidity, severe bleeding, and contraindication to anticoagulation. Precautions include careful patient selection and neuroprognostication, as outcomes depend heavily on candidacy.

\section*{Risks}
Major bleeding, circuit thrombosis, hemolysis, infection, renal failure, limb ischemia from arterial cannulation, and neurologic complications including stroke. Mechanical or pump failure is rare but catastrophic.

\section*{Aftercare and Recovery}
Patients are managed in an intensive care unit with daily circuit and end-organ assessments while underlying disease is treated. Support is weaned as native function recovers, and decannulation occurs once the patient is stable on conventional support.

\section*{Outcome and Follow-up}
Improvement is tracked by arterial blood gases, hemodynamics, and echocardiographic function of the native heart and lungs. Persistent organ failure or failure to wean portends poor outcome and prompts reassessment of goals of care.

\section*{Expected Results}
Stable oxygenation and ventilation on reduced circuit flows, adequate perfusion and mixed venous saturation, and recovery of native cardiac and respiratory function permitting decannulation.

\section*{Follow-up}
After decannulation, continued rehabilitation, imaging for sequelae of critical illness, and long-term neurocognitive and pulmonary follow-up are required. Survivors often need physical and occupational therapy.

\section*{References}
\begin{enumerate}
  \item MedlinePlus. Extracorporeal membrane oxygenation (ECMO). U.S. National Library of Medicine; 2025.
  \item Current Medical Diagnosis and Treatment. McGraw-Hill; 2025.
\end{enumerate}

\clearpage
